%! Author = jakobnieder
%! Date = 3.01.24
\section{WEIRD Participants}
\label{sec:weired-participants}

It is a long-standing issue that the vast majority of participants in research about human psychology and behavior are WEIRD participants. WEIRD is an acronym proposed in \cite{henrichWeirdest2010} that stands for Western, Educated, Industrialized, Rich, and Democratic. Statistics from 2008 show that research recruits 96\% of their samples from only 12\% of the world population. Apparently, this recruitment window is getting narrower as 64\% - 80\% of these participants are undergraduates in psychology courses \parencite{arnettNeglected2008}.

%Problematic inference that findings are universal
    Inherently, there is no problem with this habit of recruitment. However, it is crucial to consider that inferences are only valid for the population for which the sample is representative. Without justification, a generalization of the general human psyche is not necessarily valid. The Müller-Lyer illusion is an example from the visual domain that induces differences in perception between populations \parencite{müller-lyer}. Far too often, researchers thoughtlessly claim to have discovered universality in human behavior without having tested for diversity across subpopulations. \textcite{WardUndergradsEmployed1993} discovered significant differences even between undergraduates and employed American adults. This diversity can differ in magnitude and direction, depending on the research question \parencite{henrichWeirdest2010}.

%Variation possibly originates from diverse reasons:
    Adaptation to environmental conditions through behavioral plasticity and genetic processes can lead to these variations in behavior. Cultural evolution, for instance, has a fundamental impact on the environment of individuals. Firstly, invention of tools, technology, and other knowledge impacts the physiological environment. Second, the cognitive sphere is influenced by the creation of counting systems, written symbols, categories, and heuristics. Third, the social environment is designed by norms, institutions, and laws \parencite{henrichWeirdest2010}.
    %(S. 80 rechts)

    Taking this into account, it could be posited that WEIRD participants may even represent the least suitable subgroup from which to derive generalizations. In contrast to e.g. small-scale societies, WEIRD populations typically exhibits a pronounced division of labor, considerable engagement with books and television, and a childhood experience marked by minimal physical punishment yet substantial levels of praise. Crystalline families and a high degree of individualism are further characteristics \parencite{henrichWeirdest2010}.
    %(S. 80 links)

    Encapsulation and therefore behavioral separation between sub-populations can also have genetic origins as relative similarity might reflect shared genetic origin \parencite{roseShared1988}. However, this assumption was challenged by a study showing that the expression of genes still can lead to different phenotype \parencite{kimCulture2010}. In addition, natural selection is thought to favour onto-genetic adaptation, enabling humans to better adapt to their environments through non-genetic means \parencite{henrichWeirdest2010}. Overall, evolutionary theory frequently serves as the foundation for discussions on universals although it may offer arbitrary explanations for varied behaviors. 

Aim of this argument is not to prove all populations to be substantially different in all types of behavior. Indeed, in many domains, fundamental similarities can be found across all populations such as the ability to make music, learn a language and theory of mind \parencite{henrichWeirdest2010}. However, it tries to question the predominant inference to carelessly conclude from WEIRED people that represent a very narrow sample of the world's population on universals. In fact, science can only benefit from this criticism to further expand the body of evidence and provide a more comprehensive picture of human behaviour.

In the field of color cognition the world-color-survey certainly was a pioneer project. Not only the sheer size of the dataset is impressive and makes it the most influential study related to color terms. The languages examined in the WCS all came from non-industrialized cultures and unwritten languages. Other studies tried to broaden this picture for languages which are spoken by many speakers such as English \parencite{lindseyColor2014}, Spanish and Chinese \parencite{xuComparison2023} and Russian \parencite{winawerRussian2007}. However, there is still a lack of completeness in this area of languages form industrialized and globalized countries. With the help of a new crowd-sourcing approach, in which we conducted online experiments worldwide, we were able to collect data from 25 mostly industrialized countries. In this way, this research can help close this long-standing knowledge gap and complete the picture.




% Example visual perception
    %when even the effect of carpentered corners have such an strong effect, we suspect that this might even stronger for the phenomenon related to language.
    %S. 64
%Even differences within the population of american participants.
    %From there they even go ahead and show fundamental differences between the broad amarican polulation to the psychology undergraduates... and childere
    %S(76)
%Also similarities detected(S. 78)
%Evolution theory often basis on which universals discussed but this can have arbitrary explanations for different behaviour.
    %also onto-genetic adaptation suspected to be favored by natural selection to make human capable of adapting non-genetically on environment.
%criteria to anticipate universality
%Strategies to change this habit.
    %institutionalized incentives.
    %Both from research institutions and publication side.
    %setting up testing facilities at but terminals, rail stations or airports
    %we got much better solution to this!


% Such as genetic, cultural, or environmental.
% In the context of color naming, the most prominent example is the difference between the English and the Himba language.
% The Himba language has only five color terms, whereas English has eleven.
% \parencite{sturgesStudies1969}
% This difference is often explained by the fact that the Himba language has no word for blue.
% \parencite{berlinBasic1969}
% This explanation is based on the Sapir-Whorf hypothesis, which states that the language we speak influences the way we perceive the world.
% \parencite{whorfLanguage1956}
% This hypothesis is controversial and has been criticized for its lack of empirical evidence.
% \parencite{regierColor2007}
% Nevertheless, it is still a popular explanation for the differences in color naming patterns between languages.
% \parencite{regierColor2007}



%This is especially important in the context of color naming because it is often used as an example of a universal phenomenon.
%\parencite{regierColor2007}

